\scanpage{137}
expliciter \deleted{\unsure{l}}cette action en tant qu'action sur $\mathscr D$ (défini par (22)).

Si une structure complexe sur $\mathbb R^2$ est définie par l'élément $z\in\mathscr D$, $z=\deleted{\unsure{\Re}}x+iy$, alors l'application \inserted{$\mathbb C$-linéaire} $\varphi\colon L_z=\mathbb C\cdot(z,1)\longrightarrow V$ s'identifie à
\[
 \underset{\substack{\Vert\\a+ib}}{\lambda}\longmapsto
 \Re\bigl(\deleted{\unsure{\lambda(z,1)}}\underbrace{\lambda(ze_1+e_2)}_{[(ax-by)+i(ay+bx)]e_1+(a+ib)e_2}\bigr)=(ax-by)e_1+ae_2
\]
les images inverses $\lambda_1,\lambda_2$ de $e_1,e_2$ par celle-ci sont données par
\[
 \lambda_1=-\frac{i}{y}\qquad\lambda_2=1+i\frac{x}{y}
\]
d'où
\[
 \lambda_2/\lambda_1=i(y+ix)=-x+iy
\]
donc on a
\[
 \underbrace{e_2/e_1}_{\substack{\text{quotient dans }\mathbb C^*\\
 \text{au sens de la}\\\text{structure compl. de }\mathbb R^2\\
 \text{définie par }z=x+iy\in\mathscr D}}
 =(-x+iy)^{\overline{\phantom{x}}}=-x-iy=\underset{\in\mathscr D^-}{-z}
\]

Donc si une structure de Torelli sur $E$ est définie par la \deleted{\unsure{élément|éléments}}\inserted{base} $(z_1,z_2)$ de $\Pi\subset\mathfrak t_E$, alors l'élément $z\in\mathscr D$ correspondant est donné par
\[
 \deleted{z_2/z_1=-x+iy=-\bar z\quad\text{où }z=x+iy.}
 \quad\inserted{z=-(z_2/z_1)} \tag*{(24)}
\]

Soit alors
\[
 g\in\mathrm{Sl}(2,\mathbb Z)\qquad
 g=\begin{pmatrix}a&b\\c&d\end{pmatrix}\qquad
 \begin{aligned}a,b,c,d&\in\mathbb Z\\ad-bc&=1\end{aligned}
\]
et soit \deleted{\unsure{un|une} \unsure{$u$|une}}
\[
 u\colon\mathbb Z^2\xrightarrow{\sim}\Pi
\]
une structure de Torelli, correspondant à
% archive sheet number: 136
